% ============================================================
% Thesis preamble: packages, layout, theorem environments,
% draft-marking commands. Union of the preambles of the papers
% being compiled into the thesis, with conflicts resolved.
% Chapter-specific macros live in macros/<chapter>.tex.
% ============================================================

% --- layout (University of Warwick guidelines: A4, >=40mm binding
% --- margin, 1.5 line spacing for the main text) ---
\usepackage[a4paper, left=40mm, right=25mm, top=25mm, bottom=25mm]{geometry}
\usepackage{setspace}
\onehalfspacing
\usepackage{emptypage}

% --- fonts / typography ---
\usepackage[T1]{fontenc}
\usepackage[utf8]{inputenc}
\usepackage{lmodern}
\usepackage{microtype}

% --- colour ---
\usepackage[dvipsnames]{xcolor}

% --- math ---
\usepackage{amsmath}
\usepackage{amsfonts}
\usepackage{amssymb}
\usepackage{amsthm}
\usepackage{mathtools}
\usepackage{mathrsfs}
\usepackage{nicefrac}
\usepackage{xfrac}

% --- programming-style command definitions ---
\usepackage{xparse}
\usepackage{xargs}

% --- graphics, figures, tables ---
\usepackage{graphicx}
\usepackage{caption}
\usepackage{subcaption}
\usepackage{booktabs}
\usepackage{longtable}
\usepackage{tabularx}
\usepackage{threeparttable}
\usepackage{wrapfig}
\usepackage{float}
\usepackage{tikz}
\usetikzlibrary{calc}
\usepackage{pgfplots}
\pgfplotsset{compat=1.17}

% --- code and algorithms ---
\usepackage{verbatim}
\usepackage{algorithm}
\usepackage[noend]{algpseudocode}
\usepackage{pythonhighlight}

% --- text formatting ---
\usepackage{enumitem}
\usepackage{changepage}
\usepackage{mdframed}
\usepackage{pifont}
\newcommand{\xmark}{\ding{55}}
\newcommand{\cmark}{\ding{51}}

% --- boxes (promptbox etc., shared by several chapters) ---
\usepackage{tcolorbox}

% --- citations ---
\usepackage{natbib}

% --- appendices within chapters ---
\usepackage{appendix}

% --- hyperlinks (load late), cleveref (load last) ---
\usepackage[hyphens]{url}
\usepackage[
  colorlinks=true,
  linkcolor=MidnightBlue,
  urlcolor=MidnightBlue,
  citecolor=MidnightBlue,
  anchorcolor=MidnightBlue,
]{hyperref}
\usepackage[capitalize,noabbrev]{cleveref}

% ============================================================
% Draft marking.
%
% All "glue" text composed for the thesis (i.e. not lifted from
% the underlying papers/reports) is wrapped in \draft{...} or
% \begin{draftblock}...\end{draftblock}. It renders in dark red
% with a marginal [draft] tag. Flip \draftfinal to typeset
% draft text as normal text (for the final submission).
% ============================================================
\colorlet{draftcolor}{BrickRed}
\newif\ifdraftfinal
\draftfinalfalse   % \draftfinaltrue -> draft text typesets normally
\NewDocumentCommand{\draftmark}{}{%
  \ifdraftfinal\else\marginpar{\scriptsize\color{draftcolor}[draft]}\fi}
\NewDocumentCommand{\draft}{+m}{%
  \ifdraftfinal#1\else{\color{draftcolor}#1}\draftmark\fi}
\NewDocumentEnvironment{draftblock}{}
  {\ifdraftfinal\else\color{draftcolor}\fi\draftmark}
  {}

% \chapterabstract{...}: paper abstracts re-used as chapter openers.
\NewDocumentCommand{\chapterabstract}{+m}{%
  \begin{quote}\itshape #1\end{quote}\medskip}

% \chaptersource{...}: note attributing a chapter to its source paper.
\NewDocumentCommand{\chaptersource}{+m}{%
  \begin{mdframed}[backgroundcolor=black!5, linewidth=0pt]
  \small #1
  \end{mdframed}\medskip}

% ============================================================
% Theorem environments (shared by all chapters; numbered within
% chapter so that e.g. Theorem 5.3 is the third result of Ch. 5).
% ============================================================
\theoremstyle{plain}
\newtheorem{theorem}{Theorem}[chapter]
\newtheorem{nontheorem}[theorem]{Non-theorem}
\newtheorem{proposition}[theorem]{Proposition}
\newtheorem{lemma}[theorem]{Lemma}
\newtheorem{corollary}[theorem]{Corollary}
\theoremstyle{definition}
\newtheorem{definition}[theorem]{Definition}
\newtheorem{lefinition}[theorem]{Predefinition}
\newtheorem{example}[theorem]{Example}
\newtheorem{assumption}[theorem]{Assumption}
\theoremstyle{remark}
\newtheorem{remark}[theorem]{Remark}
\newtheorem*{remark*}{Remark}
\newtheorem*{notation*}{Notation}
\newtheorem*{setup*}{Setting}
\newtheorem*{notes*}{Explanatory notes}
\newtheorem*{constants*}{Constants}

% proof sketch
\newenvironment{hproof}{%
  \renewcommand{\proofname}{Proof sketch}\proof}{\endproof}

% centred bold-italic "thesis statement" display
\newenvironment{thesisclaim}{%
  \begin{center}\bfseries\itshape}{\end{center}}

% ============================================================
% Todo notes (disabled by default; commands kept so lifted text
% containing them still compiles).
% ============================================================
\usepackage[disable, textsize=tiny]{todonotes}
\newcommandx{\manyu}[2][1=]{\todo[linecolor=black,backgroundcolor=black!25,bordercolor=black,#1]{APS: #2}}
\newcommandx{\manyuunimp}[2][1=]{\todo[linecolor=black!50,backgroundcolor=black!12.5,bordercolor=black!50,textcolor=black!50,#1]{APS: #2}}
\newcommandx{\towrite}[2][1=]{\todo[linecolor=black,backgroundcolor=black,bordercolor=black,textcolor=white#1,inline]{ADD: #2}}
\newcommandx{\rephrase}[2][1=]{\todo[linecolor=black,backgroundcolor=black!25,bordercolor=black,#1]{OR: #2}}
\newcommandx{\bleg}[2][1=]{\todo[linecolor=blue,backgroundcolor=blue!25,bordercolor=blue,#1]{Q: #2}}
\newcommandx{\dpaleka}[2][1=]{\todo[linecolor=cyan,backgroundcolor=cyan!25,bordercolor=cyan,#1]{DP: #2}}
\newcommandx{\Alejandro}[2][1=]{\todo[linecolor=OliveGreen,backgroundcolor=OliveGreen!25,bordercolor=OliveGreen,#1]{Al: #2}}
\newcommandx{\vin}[2][1=]{\todo[linecolor=magenta,backgroundcolor=magenta!25,bordercolor=magenta,#1]{VB: #2}}
\newcommandx{\adam}[2][1=]{\todo[linecolor=blue,backgroundcolor=blue!25,bordercolor=blue,#1]{Adam: #2}}
\newcommandx{\evan}[2][1=]{\todo[linecolor=red,backgroundcolor=red!25,bordercolor=red,#1]{Evan: #2}}
\newcommandx{\ft}[2][1=]{\todo[linecolor=orange,backgroundcolor=orange!25,bordercolor=orange,#1]{FT: #2}}

% ============================================================
% Shared colours and boxes (from the cfcasting/solib preambles)
% ============================================================
\definecolor{pastelred}{HTML}{FFC0C0}
\definecolor{pastelorange}{HTML}{FFDAC0}
\definecolor{pastelyellow}{HTML}{FFFDC0}
\definecolor{pastelgreen}{HTML}{a0e7a0}
\definecolor{pastelcyan}{HTML}{C0FFFD}
\definecolor{pastellightblue}{HTML}{D0E0FF}
\definecolor{pastelblue}{HTML}{C0C0FF}
\definecolor{pastelviolet}{HTML}{DAC0FF}
\definecolor{pastelmagenta}{HTML}{FFC0FF}
\definecolor{pastelrose}{HTML}{FFC0DA}
\definecolor{darkred}{HTML}{C23B22}
\definecolor{darkgreen}{HTML}{1cc650}
\definecolor{lightgreen}{HTML}{caee9c}

\newenvironment{promptbox}{
  \begin{figure}[hbt!]
  \begin{tcolorbox}[left=1.5mm, right=1.5mm, top=1.5mm, bottom=1.5mm]
  \raggedright
  \small
  \textbf{User:} \texttt{...} \\[4pt]
  \textbf{Assistant:} \texttt{...} \\[4pt]
}{
  \end{tcolorbox}
  \end{figure}
}

\newenvironment{emptypromptbox}[1][]
{
  \begin{tcolorbox}[left=1.5mm, right=1.5mm, top=1.5mm, bottom=1.5mm]
    \raggedright
    \small
    \ifx\relax#1\relax\else
      \begin{center}
        {\normalsize \textbf{\color{black} #1}}
      \end{center}
    \fi
}{
  \end{tcolorbox}
}

\newtcolorbox{questionbox}[2][]{colback=gray!10!white, colframe=gray!40!black, title=#2, #1}

% misc text helpers
\newcommand{\nicett}[1]{{\texttt{#1}}}
